\section{模型假设}\label{sec:assumptions}

本节集中列出建模所用的假设。每条给出依据与失效后果；后文不再重复声明。

\begin{enumerate}
  \item \textbf{基因组先于网络结构。} 个体由二倍体基因组 $G$（离散碱基串与调控基序）描述，
        碱基不直接编码权重，只经基因调控网络影响细胞命运，网络结构是发育算子的产物。
        依据：genomic bottleneck \citep{bib1_zador2019,bib3_shuvaev2024}。
        失效后果：若碱基直接编码权重，发育层的解释力消失，结构变异与基因变异将不再可分。

  \item \textbf{二倍体结构与等位基因效应的通路单一化。} 个体含两对同源染色体，
        等位基因效应只经调控读出 $q(G)$ 传给基因调控网络。
        依据：孟德尔式结构分离需要二倍体与独立分配；通路单一化使单碱基与结构变化可归因。
        失效后果：位点表达量若另行进入发育，则变异来源混杂，单碱基扰动的效应不可辨识。

  \item \textbf{同源链聚合取加性平均，调控读出用失配等权的简化相似度。}
        依据：加性是定量遗传的标准基线 \citep{bib138_fisher1918}，简化相似度可解释、可编辑。
        失效后果：加性会抹平完全显性；简化式丢位置权重，须以 PWM 与逐基序取最大作对照（\S\ref{sec:validation}）。

  \item \textbf{发育过程有限且同分布。} 初始前体细胞均分到六个功能域，每个前体至多分裂一次，
        神经元数 $N\in[24,48]$；位置取域内均匀随机。
        依据：发育式连接生成需要「同一机制、不同随机实现」的可控变异性。
        失效后果：若分裂次数无上界，结构规模不受控，对照实验的固定容量假设失效。

  \item \textbf{连接概率由类型兼容、空间代价与表达相似度三项决定。}
        依据：线性距离项等价于指数距离规则，有连接组学支撑 \citep{bib125_ercseyravasz2013}；
        随机放置已足以复现特异功能连接 \citep{bib123_hill2012}。
        失效后果：空间代价系数有量纲、绑定归一化发育方域，若改用文献的毫米口径则不可比。

  \item \textbf{突触符号由突触前类型决定（Dale 约束），时间常数逐神经元异质。}
        依据：抑制性前神经元自抑制等定性符号模式 \citep{bib14_dorkenwald2024}；
        异质整合时间与 NMDA、GABA$_B$ 及数百毫秒决策整合同量级。
        失效后果：时间常数的语义是网络级整合常数，区别于单神经元膜常数的量级。

  \item \textbf{神经元动力学取离散漏积分发放形式，非线性取双曲正切。}
        依据：$\tau\ge1$ 时状态是自身与有界非线性的凸组合，零输入下状态有界。
        失效后果：若 $\tau<1$ 则 $1/\tau>1$ 放大状态，有界性不再成立。

  \item \textbf{感官增益、输入权重偏置与神经元偏置按细胞类型固定，不参与学习。}
        依据：保持发育输出契约 $(A,Z,\tau,W^{(0)},M)$ 冻结，感官通路仍经细胞类型与基因型挂钩。
        失效后果：若这些量可学习，则「结构由发育生成」的主张被稀释，参数量对照口径改变。

  \item \textbf{当代学习只改连接权重幅度，支撑与符号在训练期冻结。}
        依据：genomic bottleneck 要求后天学习不得遗传；结构只能在代际经基因组变异改变。
        失效后果：结构在当代不可改变，这是本文要测量的代价之一（对应 H1、H3）。

  \item \textbf{适应度取四项行为指标代内归一后的加权和，选择用二元锦标赛。}
        依据：四项覆盖生存、觅食、避敌、能量四个生态维度；锦标赛对适应度的平移与线性变换不变
        \citep{bib136_blickle1996,bib137_eiben2015}，而比例选择强度与适应度尺度不可分离 \citep{bib135_goldberg1991}。
        失效后果：$F$ 只是代内相对分数，不可跨代、跨环境直接比较；权重属设计选择，须做敏感性分析。

  \item \textbf{环境为连续二维闭环，边界反射，回合长度固定。}
        依据：反射边界确定性、不消耗随机数，跨种子对照才成立；固定长度使分母固定、跨回合可比。
        失效后果：幸存个体在全灭后仍空转至回合上限，这些步不携带行为信息（代价被接受）。

  \item \textbf{突变率 $\mu=0.001$ 每碱基每配子属计算抽象参数。}
        依据：与数字演化惯例同量级（Avida 默认约 $0.0075$）。
        失效后果：它比脊椎动物真实种系碱基替换率高约 $10^{5}$ 倍，不可作生物学参数解读。
\end{enumerate}
